\scanpage{4}
Suffisance. Grâce à \textcircled{\scriptsize i} et Lemme 1 on est ramené à voir que $F_{\widehat A}$ est représentable, car posant $B=\widehat A\times A_f$, $F_B$ sera alors représentable en vertu de \textcircled{\scriptsize iv}, donc $F$ l'est ; le $P$ le représentant sera alors de t.f. grâce à \textcircled{\scriptsize iii}.

Pour tt $n$, $F_{\underbrace{A/f^{n+1}A}_{A_n}}$ est représentable, par hyp. \inserted{\textcircled{\scriptsize{v}}} par un module $P_n$ sur $A_n$, \deleted{Soit $P$} de type fini grâce à (iii). Soit $\widehat P=\varprojlim P_n$, qui est de t.f. sur $\widehat A$ (N.B. les $P_n$ ``se recollent''). \deleted{Pour \unsure{tout module}} Utilisant \textcircled{\scriptsize ii}, on trouve que $\widehat P$ représente $F_{\widehat A}$ sur la catégorie des $A$-modules de t.f., et utilisant \textcircled{\scriptsize iii}, on trouve que $\widehat P$ représente $F_{\widehat A}$ partout.

\uline{Proposition} 3 \quad $A$ noeth. Pour que $F$ soit repr. par un module \deleted{de} t.f. il f. et s. que $F$ satisfasse aux conditions suivantes :

(i) $F$ exact à g.

(ii) $F$ commute aux $\varinjlim$ filtrantes.

% A bracket encloses the last physical condition; its arrowhead falls between (ii) and the condition on quotients.
\margin{\tikz{\draw[->] (0,0) .. controls (-.5,.6) and (-.45,1.05) .. (.12,1.16);}}
\inserted{(\deleted{iv}iii) Pour tout $f\in A$, et tout module de type fini $M$ sur $\widehat A$\inserted{$=\varprojlim A/f^{n+1}A$}, \unsure{on}
\[
F(\deleted{\widehat M}M)\xrightarrow{\sim}\varprojlim F(M/f^{n+1}M)
\]
\deleted{où $\widehat M=\varprojlim M/f^{n+1}M$.}}

(\deleted{iii}iv) Pour tout quotient \inserted{non nul} $B$ de $A$, \deleted{et tout $f\in B$,} \inserted{il existe un} $f\in B$, non nilpotent, tel que $F_{B_f}$ soit représentable.

% archive sheet number: 3
